\documentclass[11pt]{article}

\usepackage[a4paper,margin=1in]{geometry}
\usepackage{amsmath,amssymb,mathtools}
\usepackage{booktabs,array}
\usepackage{enumitem}
\usepackage[hidelinks]{hyperref}
\setlength{\emergencystretch}{3em}

\newcommand{\T}{\mathcal T_h}
\newcommand{\E}{\mathcal E_h}
\newcommand{\Eint}{\mathcal E_h^{\mathrm{int}}}
\newcommand{\Edir}{\mathcal E_h^{\mathrm{Dir}}}
\newcommand{\p}{\mathbb P}
\newcommand{\dof}{\operatorname{dof}}
\newcommand{\I}{\mathcal I}
\newcommand{\uhat}{\underline{\widehat u}}
\newcommand{\uK}{\underline u^K}
\newcommand{\qK}[1]{\underline q_{#1}^K}
\newcommand{\UK}{\underline U^K}
\newcommand{\FK}{\underline F^K}
\newcommand{\LK}{\mathbb L^K}
\newcommand{\SK}{\mathbb S^K}
\newcommand{\B}{\mathbb B}
\newcommand{\C}{\mathbb C}
\newcommand{\Mhat}{\widehat M}
\newcommand{\Mbb}{\widehat{\mathbb M}}
\newcommand{\Mhh}{\widehat{\widehat M}}
\newcommand{\kappab}{\boldsymbol\kappa}
\newcommand{\etab}{\boldsymbol\eta}
\newcommand{\inner}[2]{\left\langle #1,#2\right\rangle}
\newcommand{\code}[1]{\texttt{#1}}

\title{Diffusion--Reaction HDG: Derivation and Assembly Workflows}
\author{}
\date{}

\begin{document}
\maketitle

\begin{abstract}
This note gives a compact derivation-to-code map for the LDG-H/HDG discretization of a diffusion--reaction problem. It derives the local mixed system, the condensed trace equation, and the corresponding assembly logic for two implementations: a vectorized NumPy workflow and a fused Numba workflow. Up to the implementation specialization, the derivation allows a full, non-diagonal diffusion tensor $\kappab$. The implementation section then states the scalar-Laplacian specialization used in the current code.
\end{abstract}

\section{Continuous problem and mixed form}

Let $\Omega\subset\mathbb R^2$ be polygonal. Given $f$, $u_D$, a reaction coefficient $r$, and a uniformly positive diffusion tensor
\[
    \kappab(x)=
    \begin{bmatrix}
        \kappa_{11}(x) & \kappa_{12}(x)\\
        \kappa_{21}(x) & \kappa_{22}(x)
    \end{bmatrix},
\]
consider
\begin{equation}\label{eq:model}
    -\nabla\cdot(\kappab\nabla u)+ru=f\quad\text{in }\Omega,
    \qquad
    u=u_D\quad\text{on }\partial\Omega .
\end{equation}
Introduce the diffusive flux
\[
    q=-\kappab\nabla u .
\]
Let
\[
    \etab:=\kappab^{-1}
    =
    \begin{bmatrix}
        \eta_{11} & \eta_{12}\\
        \eta_{21} & \eta_{22}
    \end{bmatrix}.
\]
Then
\begin{equation}\label{eq:mixed-strong}
    \nabla\cdot q+ru=f,
    \qquad
    \etab q=-\nabla u.
\end{equation}
On each element $K\in\T$, integration by parts gives
\begin{align}
    \inner{q\cdot n}{v}_{\partial K}
    -\inner{q}{\nabla v}_{K}
    +\inner{ru}{v}_{K}
        &=\inner{f}{v}_{K}, \label{eq:weak-balance}\\
    -\inner{\etab q}{w}_{K}
    +\inner{u}{\nabla\cdot w}_{K}
        &=\inner{u}{w\cdot n}_{\partial K}. \label{eq:weak-flux}
\end{align}
HDG replaces the boundary value $u|_{\partial K}$ by a single-valued trace unknown $\widehat u$ on mesh edges.

\section{LDG-H formulation}

Use
\begin{align*}
    V_h &= \{v\in L^2(\Omega):v|_K\in\p_p(K),\ K\in\T\},\\
    W_h &= \{w\in L^2(\Omega)^2:w|_K\in\p_p(K)^2,\ K\in\T\},\\
    M_h &= \{\mu\in L^2(\partial\T):\mu|_\Gamma\in\p_p(\Gamma),\ \Gamma\in\E\}.
\end{align*}
The LDG-H normal flux is
\begin{equation}\label{eq:num-flux}
    \widehat q_h\cdot n=q_h\cdot n+\tau(u_h-\widehat u_h),
    \qquad \tau>0.
\end{equation}
For every $K\in\T$, the local equations are
\begin{align}
    \inner{q_h\cdot n}{v}_{\partial K}
    -\inner{q_h}{\nabla v}_{K}
    +\inner{ru_h}{v}_{K}
    +\inner{\tau u_h}{v}_{\partial K}
        &=\inner{f}{v}_{K}+\inner{\tau\widehat u_h}{v}_{\partial K},
        &&v\in\p_p(K), \label{eq:hdg-u}\\
    -\inner{\etab q_h}{w}_{K}
    +\inner{u_h}{\nabla\cdot w}_{K}
        &=\inner{\widehat u_h}{w\cdot n}_{\partial K},
        &&w\in\p_p(K)^2. \label{eq:hdg-q}
\end{align}
On Dirichlet edges,
\begin{equation}\label{eq:dirichlet-trace}
    \widehat u_h^\Gamma=\Pi_{\p_p(\Gamma)}u_D,
    \qquad \Gamma\in\Edir.
\end{equation}
On interior edges, $\widehat u_h$ is determined by flux continuity. Define the incidence set
\begin{equation}\label{eq:incidence-set}
    \I_\Gamma:=\{(K,e):\sigma(K,e)=\Gamma\},
\end{equation}
where $\sigma(K,e)$ is the global edge attached to local edge $e$ of element $K$. If $\Gamma\in\Eint$, then
\[
    \I_\Gamma=\{(K,e),(K',e')\},
    \qquad
    \sigma(K,e)=\sigma(K',e')=\Gamma .
\]
The interior transmission condition is
\begin{equation}\label{eq:flux-continuity}
    \sum_{(K,e)\in\I_\Gamma}
    \inner{q_h^K\cdot n^K+\tau^{K,e}(u_h^K-\widehat u_h^\Gamma)}{\mu}_{\Gamma}=0,
    \qquad
    \mu\in\p_p(\Gamma).
\end{equation}
This form emphasizes that the global edge equation is assembled from local element-side contributions.

\section{Element and trace matrices}

On $K$, expand
\[
    u_h^K=\sum_{j=1}^{n_K}u_j^K\varphi_j^K,
    \qquad
    (q_h^K)_\ell=\sum_{j=1}^{n_K}q_{\ell,j}^K\varphi_j^K,
    \qquad
    n_K=\dof(K)=\frac{(p+1)(p+2)}2.
\]
On $\Gamma$, expand
\[
    \widehat u_h^\Gamma=\sum_{j=1}^{n_\Gamma}\widehat u_j^\Gamma\widehat\varphi_j^\Gamma,
    \qquad n_\Gamma=\dof(\Gamma)=p+1.
\]
We use the following matrices:
\begin{align*}
    M_\alpha^K[i,j]
        &=\inner{\alpha\varphi_j^K}{\varphi_i^K}_{K}, &
    D_\ell^K[i,j]
        &=\inner{\partial_{x_\ell}\varphi_j^K}{\varphi_i^K}_{K},\\
    M_\theta^{\partial K}[i,j]
        &=\inner{\theta\varphi_j^K}{\varphi_i^K}_{\partial K}, &
    \Mhat_\theta^{\partial_eK}[i,j]
        &=\inner{\theta\widehat\varphi_j^{\sigma(K,e)}}{\varphi_i^K}_{\partial_eK},\\
    \Mhat_\theta^{\sigma(K,e),K}[i,j]
        &=\inner{\theta\varphi_j^K}{\widehat\varphi_i^{\sigma(K,e)}}_{\sigma(K,e)}, &
    \Mhh_\theta^{K,e}[i,j]
        &=\inner{\theta\widehat\varphi_j^{\sigma(K,e)}}{\widehat\varphi_i^{\sigma(K,e)}}_{\sigma(K,e)}.
\end{align*}
The tensor inverse $\etab=\kappab^{-1}$ gives four scalar coefficient mass matrices
\[
    M_{\eta_{\ell m}}^K[i,j]
    :=
    \inner{\eta_{\ell m}\varphi_j^K}{\varphi_i^K}_{K},
    \qquad
    \ell,m\in\{1,2\}.
\]
Thus no hat means element--element, one hat means trace--element or element--trace coupling, and two hats mean trace--trace products. This is why the local trace-mass contribution is denoted by $\Mhh_\theta^{K,e}$ rather than by a generic block symbol.

Testing \eqref{eq:hdg-u} with $v=\varphi_i^K$ gives
\begin{equation}\label{eq:matrix-u}
    (M_r^K+M_\tau^{\partial K})\uK
    +\sum_{\ell=1}^2
    \left(M_{n_\ell}^{\partial K}-(D_\ell^K)^T\right)\qK{\ell}
    =
    \underline f^K+
    \sum_{e=1}^3\Mhat_\tau^{\partial_eK}\uhat^{\sigma(K,e)} .
\end{equation}
Testing \eqref{eq:hdg-q} with $w=\varphi_i^K e_\ell$ gives the coupled flux-component equations
\begin{equation}\label{eq:matrix-q-component}
    (D_\ell^K)^T\uK
    -
    \sum_{m=1}^2 M_{\eta_{\ell m}}^K\qK{m}
    =
    \sum_{e=1}^3\Mhat_{n_\ell}^{\partial_eK}\uhat^{\sigma(K,e)},
    \qquad
    \ell=1,2 .
\end{equation}
Equivalently,
\begin{align}
    (D_1^K)^T\uK
    -M_{\eta_{11}}^K\qK{1}
    -M_{\eta_{12}}^K\qK{2}
        &=
    \sum_{e=1}^3\Mhat_{n_1}^{\partial_eK}\uhat^{\sigma(K,e)}, \label{eq:q1-full-tensor}\\
    (D_2^K)^T\uK
    -M_{\eta_{21}}^K\qK{1}
    -M_{\eta_{22}}^K\qK{2}
        &=
    \sum_{e=1}^3\Mhat_{n_2}^{\partial_eK}\uhat^{\sigma(K,e)}. \label{eq:q2-full-tensor}
\end{align}
The off-diagonal matrices $M_{\eta_{12}}^K$ and $M_{\eta_{21}}^K$ are precisely the terms absent in a diagonal or scalar diffusion model.

\section{Local solver and condensation}

Collect
\[
    \UK=\begin{bmatrix}\uK\\ \qK{1}\\ \qK{2}\end{bmatrix},
    \qquad
    \uhat^K=\begin{bmatrix}
        \uhat^{\sigma(K,1)}\\
        \uhat^{\sigma(K,2)}\\
        \uhat^{\sigma(K,3)}
    \end{bmatrix},
    \qquad
    \FK=\begin{bmatrix}\underline f^K\\0\\0\end{bmatrix}.
\]
Then
\begin{equation}\label{eq:local-block}
    \LK\UK=\FK+\Mbb^{\partial K}\uhat^K,
\end{equation}
where the full-tensor local matrix is
\begin{equation}\label{eq:L-block}
\LK=
\begin{bmatrix}
M_r^K+M_\tau^{\partial K}
& M_{n_1}^{\partial K}-(D_1^K)^T
& M_{n_2}^{\partial K}-(D_2^K)^T\\
(D_1^K)^T & -M_{\eta_{11}}^K & -M_{\eta_{12}}^K\\
(D_2^K)^T & -M_{\eta_{21}}^K & -M_{\eta_{22}}^K
\end{bmatrix}.
\end{equation}
The $m$-th trace-column block of $\Mbb^{\partial K}$ is
\begin{equation}\label{eq:Mbb-column}
    \Mbb_m^{\partial K}:=
    \begin{bmatrix}
        \Mhat_\tau^{\partial_mK}\\
        \Mhat_{n_1}^{\partial_mK}\\
        \Mhat_{n_2}^{\partial_mK}
    \end{bmatrix}.
\end{equation}
Let
\[
    \SK=(\LK)^{-1},
    \qquad
    \B_e^K:=
    \begin{bmatrix}
        \Mhat_\tau^{\sigma(K,e),K} &
        \Mhat_{n_1}^{\sigma(K,e),K} &
        \Mhat_{n_2}^{\sigma(K,e),K}
    \end{bmatrix}.
\]
For $\Gamma\in\Eint$ with $\sigma(K,e)=\sigma(K',e')=\Gamma$, the uncondensed matrix transmission equation is
\begin{equation}\label{eq:uncondensed-transmission-two-side}
    \left(\Mhh_\tau^{K,e}+\Mhh_\tau^{K',e'}\right)\uhat^\Gamma
    =
    \B_e^K\UK+\B_{e'}^{K'}\underline U^{K'}.
\end{equation}
Equivalently, for any interior edge,
\begin{equation}\label{eq:uncondensed-transmission-incidence}
    \sum_{(K,e)\in\I_\Gamma}\Mhh_\tau^{K,e}\uhat^\Gamma
    =
    \sum_{(K,e)\in\I_\Gamma}\B_e^K\UK.
\end{equation}
Substituting
\[
    \UK=\SK(\FK+\Mbb^{\partial K}\uhat^K)
\]
into \eqref{eq:uncondensed-transmission-incidence} gives the condensed trace equation
\begin{equation}\label{eq:condensed-general}
    \sum_{(K,e)\in\I_\Gamma}\Mhh_\tau^{K,e}\uhat^\Gamma
    -
    \sum_{(K,e)\in\I_\Gamma}\sum_{m=1}^3
    \C_{e,m}^K\uhat^{\sigma(K,m)}
    =
    \sum_{(K,e)\in\I_\Gamma}g_e^K,
    \qquad \Gamma\in\Eint,
\end{equation}
with
\begin{equation}\label{eq:C-g}
    \C_{e,m}^K:=\B_e^K\SK\Mbb_m^{\partial K},
    \qquad
    g_e^K:=\B_e^K\SK\FK.
\end{equation}
For the two-sided edge $\sigma(K,e)=\sigma(K',e')=\Gamma$, \eqref{eq:condensed-general} reads
\begin{align}\label{eq:condensed-two-side}
    &\left(\Mhh_\tau^{K,e}+\Mhh_\tau^{K',e'}\right)\uhat^\Gamma
    -\sum_{m=1}^3\C_{e,m}^K\uhat^{\sigma(K,m)}
    -\sum_{m=1}^3\C_{e',m}^{K'}\uhat^{\sigma(K',m)} \notag\\
    &\hspace{10em}=g_e^K+g_{e'}^{K'}.
\end{align}
Assembling \eqref{eq:condensed-general} over all interior edges, with Dirichlet traces fixed by \eqref{eq:dirichlet-trace}, gives the global trace system
\begin{equation}\label{eq:global-trace}
    \mathbf A\uhat=\underline b.
\end{equation}
After solving for $\uhat$, each element is reconstructed independently by \eqref{eq:local-block}.

\section{Implementation specialization}

The implementation discussed here targets the scalar Laplacian case
\[
    -\Delta u+ru=f,
    \qquad q=-\nabla u.
\]
Equivalently,
\[
    \kappab=I,
    \qquad
    \etab=I,
\]
so
\[
    M_{\eta_{11}}^K=M^K,
    \qquad
    M_{\eta_{22}}^K=M^K,
    \qquad
    M_{\eta_{12}}^K=0,
    \qquad
    M_{\eta_{21}}^K=0.
\]
Thus the full-tensor local matrix \eqref{eq:L-block} reduces to
\begin{equation}\label{eq:L-block-scalar}
\LK=
\begin{bmatrix}
M_r^K+M_\tau^{\partial K}
& M_{n_1}^{\partial K}-(D_1^K)^T
& M_{n_2}^{\partial K}-(D_2^K)^T\\
(D_1^K)^T & -M^K & 0\\
(D_2^K)^T & 0 & -M^K
\end{bmatrix}.
\end{equation}
The tensor-diffusion extension only changes the flux mass subblock in \eqref{eq:L-block}; the trace equation, incidence-based assembly, and static-condensation pattern are unchanged.

Both implementations assemble the same condensed system \eqref{eq:condensed-general}. They differ only in when local Schur blocks are materialized.

\section{Vectorized NumPy workflow}

The vectorized path forms the local Schur tensor
\begin{equation}\label{eq:trace-block-tensor}
    \code{trace\_blocks}[K,e,m,i,j]
    \equiv
    \C_{e,m}^K[i,j]
    =
    \B_e^K\SK\Mbb_m^{\partial K}[i,j].
\end{equation}
Before COO flattening, the column trace basis is oriented consistently with the global edge orientation. If $\pi_{K,m}$ is the identity or reversal permutation, the assembled block from row edge $\sigma(K,e)$ to column edge $\sigma(K,m)$ is
\[
    \C_{e,m}^K[:,\pi_{K,m}].
\]
For each interior row edge $\Gamma=\sigma(K,e)$, the Schur part contributes
\[
    \mathbf A_{\Gamma,\sigma(K,m)} \mathrel{+}= -\C_{e,m}^K,
    \qquad m=1,2,3,
\]
while the trace--trace mass contribution is
\[
    \mathbf A_{\Gamma,\Gamma}\mathrel{+}=\Mhh_\tau^{K,e}.
\]
The RHS contribution is
\[
    \underline b_\Gamma\mathrel{+}=g_e^K.
\]

\paragraph{Vectorized NumPy HDG assembly.}
Input: mesh and basis data, $f$, $r$, $u_D$, and stabilization $\tau$.
The output is the trace solution $\uhat$ and reconstructed local unknowns
$\{\UK\}$. The workflow is:
\begin{enumerate}[leftmargin=*,nosep]
\item Build topology: $\sigma(K,e)$, boundary labels, orientations, normals,
      and Jacobians.
\item Assemble all local matrices entering $\LK$, $\Mbb_m^{\partial K}$,
      and $\B_e^K$.
\item Compute or apply local inverse actions $\SK=(\LK)^{-1}$.
\item Materialize $\C_{e,m}^K=\B_e^K\SK\Mbb_m^{\partial K}$ for all
      $K,e,m$.
\item Orient column trace bases and flatten $-\C_{e,m}^K$ into COO arrays.
\item Insert $\Mhh_\tau^{K,e}$ into diagonal edge blocks and scatter-add
      $g_e^K=\B_e^K\SK\FK$ into the right-hand side.
\item Eliminate or impose Dirichlet trace degrees of freedom.
\item Solve the sparse trace system and reconstruct
      $\UK=\SK(\FK+\Mbb^{\partial K}\uhat^K)$ elementwise.
\end{enumerate}

\section{Fused Numba workflow}

The fused Numba path uses the same formulae, but computes local inverse actions inside compiled element loops. It avoids storing the full tensor \eqref{eq:trace-block-tensor}. Dirichlet trace unknowns are eliminated during assembly: if $\sigma(K,m)$ is Dirichlet, then
\[
    -\C_{e,m}^K\uhat^{\sigma(K,m)}
\]
is moved to the reduced RHS instead of being inserted into the reduced matrix.

\paragraph{Fused Numba reduced trace assembly.}
Input: mesh and reference data, projected arrays $f_h,r_h$, $u_D$, and
stabilization $\tau$. The output is $\mathbf A_{FF}$, $\underline b_F$, and
the reconstructed local unknowns. The workflow is:
\begin{enumerate}[leftmargin=*]
\item Project $u_D$ on Dirichlet edges, build free-trace maps, and allocate
      reduced COO arrays and right-hand-side scatter buffers.
\item For every $K\in\T$, map reference data, assemble $\LK$, $\FK$,
      $\Mbb_m^{\partial K}$, $\B_e^K$, and $\Mhh_\tau^{K,e}$, then factor
      the small dense matrix $\LK$.
\item For each free row edge $e$, add $g_e^K=\B_e^K\SK\FK$ and
      $\Mhh_\tau^{K,e}$. For each column edge $m$, compute
      $\C_{e,m}^K=\B_e^K\SK\Mbb_m^{\partial K}$.
\item Insert $-\C_{e,m}^K$ for a free column. For a Dirichlet column, add
      $+\C_{e,m}^K\uhat_D^{\sigma(K,m)}$ to the reduced right-hand side.
\item Accumulate scatter buffers, solve
      $\mathbf A_{FF}\uhat_F=\underline b_F$, and expand the known traces.
\item Reconstruct $\UK=\SK(\FK+\Mbb^{\partial K}\uhat^K)$ in a second
      compiled element loop.
\end{enumerate}

\section{Workflow comparison}

\begin{center}
\renewcommand{\arraystretch}{1.25}
\begin{tabular}{>{\raggedright\arraybackslash}p{0.22\textwidth} >{\raggedright\arraybackslash}p{0.34\textwidth} >{\raggedright\arraybackslash}p{0.34\textwidth}}
\toprule
Aspect & Vectorized NumPy & Fused Numba \\
\midrule
Schur blocks & Materializes $\C_{e,m}^K$ as a dense tensor over elements and local edges. & Computes $\C_{e,m}^K$ locally and writes COO entries directly. \\
Memory use & Higher, but easier to inspect. & Lower, suitable for larger meshes. \\
Loop structure & Mostly hidden inside broadcasting and flattening. & Explicit element/edge loops, but compiled. \\
Boundary handling & Can impose rows or eliminate known traces. & Naturally uses reduced free-trace assembly. \\
Orientation & Applied to materialized column blocks before flattening. & Applied locally before insertion or trace gathering. \\
Best use & Prototyping and debugging. & Production runs and projected-coefficient assembly. \\
\bottomrule
\end{tabular}
\end{center}

\paragraph{Execution characteristics.}
The vectorized NumPy path materializes element batches of local inverse blocks,
element-boundary blocks, trace-lifted Schur blocks, orientation-adjusted copies,
COO indices and data, and RHS work arrays. This representation is easy to
inspect, but its temporary storage grows with the number of elements and local
polynomial dimension.

The fused Numba path instead forms the same small dense element and face
operations inside compiled loops, applies orientation locally, writes reduced
COO entries directly, and moves known Dirichlet-column contributions to the
RHS immediately. The two workflows implement the same condensed operator but
have different allocation and memory-traffic profiles. Compilation warmup must
therefore be reported separately in any timing comparison.

Dated host timing and peak-memory measurements are retained in
\path{docs/research/solver_studies/diffusion_assembly_2026_07.md}; they are not
algorithm invariants.

\section{Code correspondence}

\begin{center}
\renewcommand{\arraystretch}{1.2}
\begin{tabular}{>{\raggedright\arraybackslash}p{0.43\textwidth} >{\raggedright\arraybackslash}p{0.47\textwidth}}
\toprule
Mathematical object or operation & Representative implementation name \\
\midrule
Normalize local stabilization $\tau^{K,e}$ & \code{\_normalize\_tau}; fused Numba path uses \code{\_normalize\_diffusion\_stabilization} \\
Build the block source vector $\FK$ & \code{block\_source\_moments} \\
Build local factors entering $\LK$ & \code{\_local\_solver\_pre\_mats} \\
Build the local inverse $\SK$ & \code{\_local\_solver\_scalar\_inverse}, \code{\_local\_solver\_blocks\_numpy} \\
Build $\Mbb_m^{\partial K}$ & \code{diffusion\_element\_boundary\_mats} \\
Build the trace lift $\B_e^K$ & \code{diffusion\_trace\_lift} \\
Build $\C_{e,m}^K=\B_e^K\SK\Mbb_m^{\partial K}$ & \code{element\_to\_trace\_matrix\_from\_lift} \\
Build interior trace mass blocks $\Mhh_\tau^{K,e}$ & \code{interior\_stabilization\_mass\_blocks} \\
COO index and data flattening, NumPy path & \code{trace\_matrix\_indices}, \code{trace\_matrix\_data} \\
RHS assembly, NumPy path & \code{trace\_rhs\_from\_lift} \\
Dirichlet trace projection & \code{boundary\_trace\_coefficients} \\
Trace orientation during Schur/reconstruction & \code{element\_to\_trace\_matrix\_from\_lift}, \code{element\_traces} \\
Reduced fused assembly, Numba path & \code{assemble\_projected\_diffusion\_trace\_system\_eliminated\_numba} \\
Known trace elimination & \code{eliminate\_known\_dofs} \\
Local reconstruction & \code{reconstruct\_local\_unknowns}, \code{reconstruct\_projected\_diffusion\_local\_unknowns\_numba} \\
\bottomrule
\end{tabular}
\end{center}

\section{Minimal end-to-end flow}

\begin{enumerate}[leftmargin=2em]
    \item Build topology, edge orientations, boundary labels, normals, and geometric maps.
    \item Build reference volume and face quadrature data.
    \item Assemble $\LK$, $\Mbb_m^{\partial K}$, $\B_e^K$, $\Mhh_\tau^{K,e}$, and $\FK$.
    \item Condense the cell unknowns through $\SK=(\LK)^{-1}$.
    \item Assemble \eqref{eq:condensed-general}, eliminating Dirichlet trace degrees of freedom if desired.
    \item Solve the sparse trace system.
    \item Reconstruct $(u_h,q_h)$ independently on each element.
\end{enumerate}

\paragraph{Invariant.}
The NumPy and Numba paths assemble the same condensed HDG trace system. The difference is execution strategy: the NumPy path materializes Schur tensors, while the Numba path fuses local inverse actions and COO insertion.

\end{document}
